%\chapter{ XB: Electrical Vibrations} \section{235 i}
235 i. The introduction of the telephone into practical use, and the numerous
applications to scientific experiment of which it admits, bring the subject of
alternating electric currents within the scope of Acoustics, and impose upon us
the obligation of shewing how the general principles explained in this work
may best be brought to bear upon the problems presenting themselves.  Indeed
electricity affords such excellent illustrations that the temptation to use
some of them has already (§§ 7S, 92 a, 111 b) proved irresistible. It will be
necessary, however, to take for granted a knowledge of elementary electrical
theory, and to abstain for the most part from pursuing the subject in its
application to vibrations of enormously high frequency, such as have in recent
years acquired so much importance from the researches initiated by Lodge and by
Hertz. In the writings these physicists and in the works of
Prof. J. J. Thomson \footnote{Researches in Electricity and Magnetism, 1893.},
and of Mr O. Heaviside \footnote{Electrical Papers, 1892.}
the reader will find the necessary information on that branch of the subject.

The general idea of including electrical phenomena under those of ordinary
mechanics is exemplified in the early writings of Lord Kelvin; and in his
\textit{Dynamical Theory of the Electro-magnetic Field}
\footnote{Phil. Trans. vol. 155, p. 459, 1865; Collected Works, vol. 1, p. 526.},
Maxwell gave a systematic exposition of the subject from this point of view.

\section{235 j}
235 j. We commence with the consideration of a simple electrical circuit,
consisting of an electro-magnet whose terminals are connected with the poles of
a condenser, or leyden
\footnote{This term has been approved by Lord Kelvin, On a New Form of Air
Leyden, Proc. Roy. Soc., vol. 52, p. 6, 1892},
of capacity C.  The electro-magnet may be a simple coil
of insulated wire, of resistance i, and of self-induction or inductance L. If
there be an iron core, it is necessary to suppose that the metal is divided so
as to avoid the interference of internal induced currents, and further that the
whole change of magnetism is small \footnote{Phil Mag., vol. 23, p. 225, 1887.}
Otherwise the behaviour of the iron is complicated with hysteresis, and its
effect cannot be represented as a simple augmentation of Z. : Also for the
present we will ignore the hysteresis exhibited by many kinds of leydens.

If x denote the charge of the leyden at time t, i is the current, and if -Eieos
be the imposed electro-motive force, the equation is

Lx - - RA xfC  Ift (1).

The solution of (1) gives the theory of forced electrical vibrations; but we
may commence with the consideration of the free vibrations corresponding to i =
0. This problem has already been treated in § 45, from which it appears that
the currents are oscillatory, if

R2s (LiC) (2).

The fact that the discharges of leydens are often oscillatory was suspected by
Henry and by v. Helmholtz, but the mathematical theory is due to Kelvin
\footnote{On Transient Electric Currents, Phil. Mag.  June, 1853.}.
When R is ranch smaller than the critical value in (2). a large number of
vibrations occur without much loss of amplitude, and the period t is given by

T = 27r f(CL) (3).

In (2), (3) the data may bo supposed to be expressed in C.G.S.,
electro-magnetic measure. If we introduce practical units, so that L , R C'
represent the inductance, resistance and capacity reckoned respectively in
earth-quadrants or Henrys, ohms, and microfarads
\footnote{Ohm = 10, henry = 10 , microfarad =10}, we have in place of (2)
\[ R' < 2000 \sqrt{ L'/C' } \]
(2')  and in place of (3)
T = 27r-S (C' :') (;V).
With ordinary appliances the value of t is very small; but by
including a considerable coil of insulated wire in the
discharging circuit of a leyden composed of numerous glass plates Lodge
\footnote{Proc. Roy. Inst., March 1889.} has
succeeded in exhibiting oscillatory sparks of periods as long as sh second.

If the leyden be of infinite capacity or what comes to the same thing, if it
be short-circuited, the equation of free motion reduces to

Lr Rx = 0 (4).

whence \footnote{Helmholtz, Pogg. Ann., LXXX IX., p. 505, 1851.}
\[\dot{x}=\dot{x}_o e^{ -R t/L}, \]
(5) \(\dot{x}_o\) representing the value of \(\dot{x}\)
x when  \(t= 0\). The quantity fJR is
sometimes called the time-constant of the circuit, being the time during which
free currents fall off in the ratio of e : 1.

Returning to equation (1), we see that the problem falls under the general head
of vibrations of one degree of freedom, discussed in § 46. In the notation
there adopted, i)- = (CL)'~ k R'L, E=sE IL  and the solution is expressed by
equations (4) and (5).  It is unnecessary to repeat at length the discussion
already given, but it may be well to call attention to the case of resonance,
where the natural pitch of the electrical vibrator coincides with that of the
imposed force.  The first and third terms

then (§ 46) compensate one another, and the equation reduces to Rx = E co9jyt
(6).

In general, if the leyden be short-circuited (C=x),

E

(7); so that, if p much exceed i i, the current is greatly reduced by
self-induction. In such a case the introduction of a leyden of suitable
capacity, by which the self-induction is compensated, results in a large
augmentation of current\footnote{Maxwell, Experiment in Magneto-Electric
Induction, Phil. Mug., May.  1868.}. The imposed electro
motive force may be obtained from a coil forming part of the circuit
and revolving in a magnetic field.




In any circuit, where vibrations, whether forced or free, proportional to
cos pt are in progress, we have x - prx, and thus the terms due to
self-induction and to the leyden enter into the equation in the same manner.
The law is more readily expressed if we use the stiffness p, equal to 1 C,
rather than the capacity itself. We may say that a stiffness p compensates an
inductance Zr, if  = and that an additional inductance AZ is compensated by an
additional stiffness A a, provided the above proportionality hold good. This
remark allows us to simplify our equations by omitting in the first instance
the stiffness of leydens. When the solution has been obtained, we may at any
time generalise it by the introduction, in place of Z, of Z - or Z- In

following this course we must be prepared to admit Negative values of Z.

\section{235 k Secondary circuit}
235 k. We will next suppose that there are two independent circuits with coefficients
of self-induction Z, N, and of mutual induction M, and examine what will be the
effect in the second circuit of the instantaneous establishment and
subsequent maintenance of a current x in the first circuit. At the first moment
the question is one of the function T only, where

T = iZi + i¥i;y + Vy (1);

and by Kelvin's rule (§ 79) the solution is to be obtained by making (1) a
minimum under the condition that x has the given value. Thus initially

and accordingly (§ 235 j) after time t

y = - iesmt (3).

if S be the resistance of the circuit. The whole induced current, as measured
by a ballistic galvanometer, is given by in which N does not appeal'. The
current in the secondary circuit due to the cessation of a previously
established steady current x in the primary circuit is the opposite of the
above.

A curious property of the initial induced current is at once evident from
Kelvin's theorem, or from equation (2). It appears
that, if M be given, the initial current, is greatest when N is least.
Further, if the secondary circuit consist mainly of a coil of n turns, the
initial current increases with diminishing n, F(r), although Mzcn, JVxn; and
thus In fact the small current
flowing through the more numerous convolutions has the same effect as regards
the energy of the field as the larger current in the fewer convolutions. This
peculiar dependence upon n cannot be investigated by the galvanometer, at least
without commutators capable of separating one part of the induced current from
the rest; for, as we see from (4) the galvanometer reading is affected in the
reverse direction. It is possible however to render evident the increased
initial current due to a diminished n by observing the magnetizing effect upon
steel needles. The magnetization depends mainly upon the initial maximum value
of the current, and in a less degree, or scarcely at all, upon its subsequent
duration \footnote{Phil. Mag., vol. 38, p. 1, 1869; vol.  39, p. 426. 1870}.

The general equations for two detached circuits, influencing one
another only by induction, may be obtained in the usual manner from (1) and

F = + (5).  Thus Lx My + Ri = X)

Mx Ny- Sy=Y

These equations, in a more general form, are considered in § 116. If a harmonic
force X = act in the first circuit, and the second circuit be free from imposed
force ( F = 0), we have on elimination of y shewing that the reaction of 'the
secondary circuit upon the first is to reduce the inductance by
\footnote{Maxwell, Phil Trans., voL 155, p. 459. 1865, where however \(M^2\) is
misprinted M.}

, M N .... .  p
and to increase the resistance by (9)

The formulae (8) and (9) may be applied to deal with a more general problem of
considerable interest, which arises when (as in some of Henry’s experiments)
the secondary circuit acts upon a third, this upon a fourth, and so on, the
only condition being that there must be no mutual induction except between
immediate neighbours in the series. For the sake of distinctness we will limit
ourselves to four circuits.

In the fourth circuit the current is due \textit{ex hypothesi} only to induction from
the third. Its reaction upon the third, for the rate of vibration under
contemplation, is given at once by (8) and (9) ; and if we use the complete
values applicable to the third circuit under these conditions, we may
thenceforth ignore the fourth circuit. In like manner we can now deduce the
reaction, upon the secondary, giving the effective resistance and inductance of
that circuit under the influence of the third and fourth circuits ; and then,
by another step of the same kind, we may arrive at the values applicable to the
primary circuit, under the influence of all the others. The process is
evidently general; and we know by the theorem of § 111 b that, however extended
the train of circuits, the influence of the others upon the first must be to
increase its effective resistance and diminish its effective inertia, in
greater and greater degree as the frequency of vibration increases.

In the limit, when the frequency increases indefinitely, the distribution of
currents is determined by the induction-coefficients, irrespective of
resistance, and, as we shall see presently, it is of such a character that the
currents are alternately opposite in sign as we pass along the series.

\section{235 l Initial conditions}
235 l. Whatever may be the number of independent currents, or degrees of
freedom, the general equations are always of the kind already discussed §§ 82,
103, 104, viz.

ddT dF dV dt dx

 (1) where T, F, V are (§ 82) homogeneous quadratic functions. In (1)
the coordinates denote the whole quantity of electricity
which has passed at time t, the currents being
i,, Xm etc.
When \(V =0\), it is simpler to express the phenomena by means of the currents.
Thus, in the problem of steady electric flow where all the quantities \(X\),
representing electro-motive forces, are constant, the currents are determined
directly by the linear equations

dF/dx = X1 , dF/dx2 = X2,  etc (2).

On the other hand when the question under consideration is one of initial
impulsive effects, or of forced vibrations of exceedingly high frequency,
everything depends upon T, and the equations reduce to

d IT d dT  y dt d x r

(3).  As an example we may consider the problem, touched upon at
the close of § 235 k, of a train of circuits where the mutual induction is
confined to immediate neighbours, so that
\footnote{The dots are omitted as unnecessary.}

(4) coefficients such as aj3, O14, not appearing.  If \(x_1\)
be given, either as a current suddenly developed and
afterwards maintained constant, or as a harmonic time function of high
frequency, while no external forces operate in the other circuits, the problem
is to determine
\(x_2\),
\(x_3\),
etc. so as to make \(T\) as small as possible, § 79. The
equations are easily written down, but the conclusion
aimed at is perhaps arrived at more instructively by consideration
of the function \(T\) itself. For, \(T\)
being homogeneous in
\(x_1\),
\(x_2\), etc.,
we have identically

(5) And, since when \(T\) is a minimum, dT/dxj, dT/dx , etc., all vanish,
\[ 2 T min = Xj\]
But if
\(x_2\), \(x_3\), etc., had all been zero, 2T would have been equal to
\( a_{11} x_1^2\). It is clear therefore that
\( a_{12} x_1 x_2\) is negative; or, as
\( a_{12} \) is taken positive, the sign
of \(x_2\) is the opposite to that of \(x_1\).

Again supposing r,, x both given, we must, when T is a minimum, have dTjdx y
dTjdx y etc., equal to zero, and thus As before, 2r might have been

a -H 2a, ojr, j:  -4 -O s jt Xr,

simply. The minimum value is necessarily less than this, and accordingly the
signs of and are opposite. This argument may be continued, and it shews that,
however long the series may be, the induced currents are alternately opposite
in sign \footnote{Phil Mag., vol. 38, p. 13, 1869}, a result in harmony with
the magnetizations observed by Henry.

In certain cases the minimum value of T may be very nearly zero. This happens
when the coils which exercise a mutual inductive influence are so close
throughout their entire lengths that they can produce approximately opposite
magnetic forces at all points of space. Suppose, for example, that there are
two similar coils A and jB, each wound with a double wire (-Bi, 5a), and
combined so that the primary circuit consists' of A , the secondary of A. and
joined by inductionless leads, and the tertiary of B simply closed upon itself.
It is evident that T is made approximately zero by taking = and 5:3 = - arsr jj.
The argument may be extended to a train of such coils, however long, and
also to cases where the number of convolutions in mutually reacting coils is
not the same.

In a large class of problems, where leyden effects do not occur sensibly, the
course of events is determined by T and F simply.  These functions may then be
reduced to sums of squares ; and the typical equation takes the form

ax + bx X (6).

If = 0, that is if there be no imposed electro motive forces, the solution is

x xoe- i (7).

Thus any system of initial currents flowing whether in detached or connected
linear conductors, or in solid conducting masses, may be resolved into normal
components, each of which dies down exponentially at its own proper rate.

A general property of the persistences, equal to a 6, is proved in § 92 a.
For example, any increase in permeability, due to the introduction of iron
(regarded as non-conducting), or any diminution of resistance, however local,
will in general bring about a rise in the values of all the persistences
\footnote{Brit. Assoc. Report, 1885, p. 911.}.

In view of the discussions of Chapter V it is not necessary to dwell upon the
solution of equations (1) when X is retained. The reciprocal theorem of § 109
has many interesting electrical applications; but, after what has there been
said, their deduction will present no difficulty.

\section{235 m. Contiguous wires}.
235 m. In § 111 b one
application of the general formulae to an electrical system has already been
given. As another example, also relating to the case of two degrees of freedom,
w'o may take the problem of two conductors in parallel. It is not necessary to
include the influence of the leads outside the points of bifurcation ; for
provided that there be no mutual induction between these parts and the
remainder, their inductance and resistance enter into the result by simple
addition.

Under the sole operation of resistance, the total count would divide itself
between the two conductors (of resistance's R and 8) in the parts

8  R

B + S ' if+S

and we may conveniently so choose the second coordinate that the currents in
the two conductors are in general still representing the total current in the
leads. The dissipation-function, found by multiplying the squares of the above
currents by R, 8 respectively, is

+ + (1).

Also, L, M, N being the induction coefficients of the two branches.
\[ LS' + 2MRS- NR , (L-M)S+(M-N)R R + S ,a;, + (Z-23f+iV),\]
(2).  Thus, in the notation of § 111 h,
LS'+iMR8 + NR (L-M)8 + (M-N)R a,,  + 5 j,  Jj + S

OL  Z - 2Jlf + N b n RS

R + 8' 60, bn-‘R + 8.  Accordingly by (5), (8) § 111 b,

RS  ((Z,-J7)g + (J -.V)it)

~ R + S R + S(R + Sy‘ + p‘(L-2M + y''

r,LS‘ + 2MRS + NR  (L - M) S + (M- N) R

 + sF  (R'+syiL-m+N)

((L-M)8 + (M-N)Ry (L- 2M + m (R + Sy +i (L-2M+ Ny

These are respectively the effective resistance and the effective inductance of
the combination . It is to be remarked that (L - 2Jf + Jlf) is necessarily
positive, representing twice the kinetic energy of the system when the currents
in the two conductors are + 1 and - 1.  The expressions for R' and L' may be
put into a form which for many purposes is more convenient, by combining the
component fractional terms. Thus

 RS(R + S) + p‘R (M - Ny + S(L- iff) )

 ~(R + Sy + f(L -2M + Ny '

j,  L8 + 2MRS + NR + p‘ (ZiV- il ) (L - 2M + N) , ,

   (R- sy p‘(L~2M+Ny

in which (LN - J ) is positive by virtue of the nature of T.

Ab p increases from zero, we know by the general theorem § 111 b, or from the
particular expressions (3), (4), that R continually increases and that X'
continually decreases.  When p is very small, R' RS

R + S' ,1 - + 2MRS + NR 'jR sy In this case the distribution of the main
current between the conductors is determined by the resistances, and (111 b)
the values of Rf and L' coincide respectively with 2F j , 2Tlxi. The resistance
is manifestly the same as if the currents were steady.  On the other hand, when
p is very great, R(M-Ny + S(L~My (L-2M + Ny LN-M

 Phil. Mag., vol. 21, p. 377, J. J. Thomson, op. cit. § 421

In this case the distribution of currents is independent of the resistances,
being determined in accordance with Kelvin’s theorem in such a manner that the
ratio of the currents in the two conductors is (N - M) : (L - M), As when p is
small, the values in (6) coincide with \( 2 F x^2 2 T x^2 \).

When the two wires connecting the conductors in parallel are wound closely
together, the energy of the field under high frequency may be very small.
There is an interesting distinction to be noted here dependent upon the manner
in which the connections are made.  Consider, for example, the case of a bundle
of five contiguous wires wound into a coil, of which three wires, connected in
series so as to have maximum inductance, constitute one of the branches in
parallel, and the other two, connected similarly in series, constitute the
other branch. There is still an alternative as to the manner of connection of
the two branches.  If steady currents would circulate opposite ways
(j negative), the total current is divided into two parts in the ratio 3 : 2, in
such a manner that the more powerful current in the double wire nearly
neutralises at external points the magnetic effects of the less powerful
current in the triple wire, and the total energy of the system is very small.
But now suppose that the connections are such that steady currents would
circulate the same way in both branches (M positive). It is evident that the
condition of minimum energy cannot be satisfied when the currents are in the
same direction, but requires that the smaller current in the triple wire should
be in the opposite direction to that 6f the larger current in the double wire.
In fact the currents must be as
3 to - 2 ; so that
(since on the same scale the total current is unity) the component currents in
the branches are both numerically greater than the total current which is
algebraically divided between them. And this peculiar feature becomes more and
more strongly marked the nearer
L and N approach to equality
\footnote{Phil. Mag., vol. 21, p. 376, 1886.}.

The unusual development of currents in the branches is, of course, attended by
an augmented effective resistance. In the limiting case when the
\(m\)
convolutions of one branch are supposed to coincide geometrically with one
another and with the
\(n\)
convolutions of the second branch, we have

If : Af : AT = m : Jnn : n'

and from (6) R = 0, (7)
an expression which increases without limit, as m and n approach to equality.


The fact that under certain conditions the currents in both branches of a
divided circuit may exceed the current in the mains has been verified by direct
experiment'. Each of the three currents to be compared traversed short lengths
of similar Germansilver wire, and the test consisted in finding what lengths of
this wire it was necessary in the various cases to include between the
terminals of a high resistance telephone in order to obtain sounds of equal
intensity. The variable currents were derived, from a battery and scraping
contact apparatus (§ 235 r), directly included in the main circuit.  The general
formula (3'), (4') undergo simplification when the conductors in parallel
exercise no mutual induction. Thus, when Jlf = 0,

 RSJR S) (RN jh (R Sy + p (L Ny


  (R' sy p lL '+ Ny If further K - 0, (8) and (9) reduce to S(R(R + S)+p D (R
SyVp D ' ~(R Sy p L (10).  The peculiar features of the combination are brought
out most strongly when S, the resistance of the inductionless component, is
great in comparison with R. In that case if the current be steady or slowly
vibrating, it flows mainly through R, while the resistance and inductance of
the combination approximate to R and L respectively ; but if on the other hand
the current be a rapidly vibrating one, it flows mainly through S, so that the
resistance of the combination approximates to 5, and the inductance to zero.
These conclusions are in agreement with (10).  If the branches in parallel be
simple electro-magnets, L and N are necessarily positive, and the numerator in
(9) is incapable vf vanishing. But, as we have seen, when leydens are admitted,
this restriction may be removed. An interesting case arises when the second
branch is inductionless, and is interrupted by a leyden of V PAtZ. Mag,, vol.  22,
p 495, 1686. capacity 7, so that JV' = - while at the same time R S,

The latter condition reduces the numerator in (9) to (L X)(R p- LX.

Thus L vanishes, (i) when LCp  1, and (ii) when CR L, The first alternative is
the condition that the loop circuit, considered by itself, should be
isochronous with the imposed vibrations.  The second expresses the equality of
the time-constants of the two branches. If they be equal, the combination
behaves like a simple resistance, whatever be the character of the imposed
electromotive force
\footnote{Chrystal, On the Differential Telephone, Edin Trans,vol 29, p 615, 1880}.

\section{235 n. Conductors in parallel}.
235 n. When there are more than two conductors in parallel,
the general expressions for the equivalent resistance and inductance
of the combination would be very complicated; but a few
particular cases are worthy of notice.

The first of these occurs when there is no mutual induction between the
members. If the quantities relating to the various branches be distinguished by
the suffixes 1, 2, 3, and if \(E\) be the difference of potentials at the common
terminals, we have

E = + i2i) ic, = (ip 1 2 4- iBs) ‘ta ~ (1 );

so that

y  V - - . I ipL' (2),

by which R' and L' are determined. Thus, if we write

if

we have from (2)

+ 2‘iP

Equations (3) and (4) contain the solution of the problem
\footnote{Phil. Mag., vol. 21, p. 379, 1886.}.
When  p= 0,
R'=.  1 t LRr

When on the other hand p is very great,

L' = (6).

Even when the mutual induction between
various members cannot be neglected, tolerably simple expressions can be found
for the equivalent resistance and inductance in the extreme cases of p
infinitely small or infinitely large. As has already been proved, (§111 b), the
above-mentioned quantities then coincide in value with and 2  (a, -har4+ ...) ,
and the calculation

of these values is easy, inasmuch as the distribution of currents among the
branches is determined in the first case entirely by F and in the second case
entirely by 21 Thus, when p is infinitely small, is a minimum, and the currents
are in proportion to the conductances of the several branches. Accordingly, if
the induction coefficients of the branches be denoted, as in § 111 t, by Ok,
fti2i ttisi -I and the resistances by -Rj, Jij, etc., we have

ii,(l i2,)=--JJ,(l JU=-h... 1

(l i2, + l R,-h...)‘ l if,4l R.,+ ...'

Tf  wl  + 22 - -  +  .  + 2ayJRiR2 4 2aiJRiR- 4 . . .

i sr4i7R74W+:y

A similar method applies when p = x, but the result is less simple on account
of the complication in the ratios of currents due to mutual induction
\footnote{J. J. Thornton, loc. cit. § 422}.

\section{235 o}
235 o. The induction-balance, originally contrived by Dove for use with the
galvanometer, has in recent years been adapted to the telephone by Hughes
\footnote{Phil. Mag. vol VIII, p. 50, 1879},
who has described experiments illustrating the marvellous sensibility of the
arrangement. The essential features are a primary, or battery, circuit, in
which circulates a current rendered intermittent by a make and break
interrupter, or by a simple scraping contact, and a secondary circuit containing
a telephone. By suitable adjustments the two circuits are rendered conjugate,
that is to say the coefficient of mutual induction is caused to vanish, so as
to reduce the telephone to silence. The introduction into the neighbourhood of
a third circuit, whether composed of a coil of wire, or of a simple conducting
mass, such as a coin, will then in general cause a revival of sound.

The destruction of the mutual induction in the case of two flat coils can be
arrived at by placing them at a short distance apart,
in parallel planes, and with accurately adjusted overlapping. But
in Hughes’ apparatus the
balance is obtained more symmetrically by the method of duplication. Four
similar coils are employed.  Of these two i4i, jI,, mounted at some distance
apart with their planes horizontal, and connected in series, constitute the
primary induction coil. The secondary induction coil consists in like manner
of placed symmetrically at short distances from
A1, A2, and also connected in series, but in such a manner that the induction
between A1 and B1 tends to balance the induction between A2 and B2. If the
four coils were perfectly similar, balance would be obtained when the
distances were equal. This of course is not to be depended upon, but by a
screw motion the distance between one pair, e.g. A1 and B1, is rendered
adjustable, and in this way a balance between the two inductions is obtained.
Wooden cups, fitting into the coils, are provided in such situations that a
coin resting in one of them is situated symmetrically between the
corresponding primary and secondary coils. The balance, previously adjusted, is
of course upset by the introduction of a coin upon one side, but if a perfectly
similar coin be introduced upon the other side also, balance may be restored.
Hughes found that very minute differences between coins could be rendered
evident by outstanding sound in the telephone.

The theory of this apparatus, when the primary currents are harmonic, is simple
especially if we suppose that the primary current a is given. If iTj,... be
the currents; 6,, fco, ... the resistances; Un,  ., Oi-, the inductances, the
equations for the case of three circuits are U y ji a + O sa' A'a -f =  j

We now assume that rj, a ., etc. are proportional to where j 27r is the
frequency of vibration. Thus,
ip (a-joA.. + + hitx:.: = - ipivi(  ,
ip (C, -I- + h-sX-i =  ip n- i 1
whence by elimination of a-i
X 2 Upa i + + p-ttg
ipau + J - ipUtiX, p a,ias,a, ijMa + tfi (2).  ‘

Brit. Amoc. Itep. p. 472.

From this it appears that a want of balance depending on
Oi, cannot compensate for the
action of the third circuit, so as to produce silence in the secondary circuit,
unless \(b\)
\todo{inline}[index of b is smudged, 2, 3]
be negligible in comparison with
pagt, that is unless the
time-constant of the third circuit be very great in comparison with the period
of the vibration. Otherwise the effects are in different phases, and therefore
incapable of balancing.

We will now introduce a fourth circuit, and suppose that the primary and
secondary circuits are accurately conjugate, so that
\(a_{12}=0\), and also that the mutual induction \(a_{34}=0\) between the third
and fourth circuits may be neglected.  Then

whence

(3).  It appears that
two conditions must be satisfied in order to secure a balance, since both the
phases and the intensities of the separate effects must be the same. The first
condition requires that the time-constants of the third and fourth circuits be
equal, unless indeed both be very great, or both be very small, in comparison
with the period. If this condition be satisfied, balance ensues when (4); and
it is especially to be noted that the adjustment is independent of pitch, so
that (by Fourier's theorem) it suffices whatever be the nature of the variable
currents operative in the primary.

As regards the position of the third and fourth circuits, usually represented
by coins in illustrative experiments, it will be seen from the symmetry of the
right-hand member of (3) that the middle position between the primary and
secondary coils is suitable, inasmuch as the product OisOg is stationary in
value when the coin is moved slightly so as to be nearer say to the primary
and further from the secondary
\footnote{Lodge, Phil Mag., vol. 9, p. 123, 1880}.
Approximate independent of other displacements
is secured by the geometrical symmetry of the coils round the axis.

\section{235 p Approximate Balance}
235 p. For the accurate comparison of electrical quantities the '‘bridge''
arrangement of Wheatstone is usually the most convenient, and is equally
available with the galvanometer in case of steady or transitory currents, or
with the telephone in to increase of periodic currents. Similar effects may be
obtained in most cases without a bridge by the employment of the differential
galvanometer or the differential telephone \footnote{Chrystal, Edin, Trans loc cit}

In the ordinary use of the bridge the four membranes a. b. c, d combined in a
quadrilateral Fig. (53 a) are simple resistances. The battery branch f joins
one pair of opposite corners, and the indicating instrument is in the '‘bridge”
e joining the other pair.  Balance is obtained, when ad = be. But for our
purpose we have to suppose that any member, e.g. a, is not merely a resistance,
or even a combination of resistances. It may include an electromagnet, and it
may be interrupted by a leyden. But in any case, so long as the current x is
strictly harmonic, proportional to the general relation between it and the
difference of potentials V at the extremities is given by

F=(a,--ia,.)a; (1),
where and ta.j are the real and imaginary parts of a complex coefficient a, and
are functions of the frequency pl2ir. In the particular case of a simple
conductor, endowed with inductance L, Ui represents the resistance, and a is
equal to pL. In general, Uj is positive; but a may be either positive, as in
the above example, or negative. The latter case arises when a resistance R is
interrupted by a leyden of capacity C. Here fti = li, a.j = - I pC.  If there
be also inductance L,

fli = R, a. - pL - IjpC

(2).  As we have already seen, § 235 j, a2 may vanish
for a particular frequency, and the combination is then equivalent to a simple
resistance. But a variation of frequency gives rise to a positive or
negative a2.

In all electrical problems, where there is no mutual induction,
the generalized quantities, a, 6, etc., combine, just as they do when they
represent simple resistances
\footnote{
For a more complete discussion of this subject
see Heaviside On Resistance and Conductance Operators, Phil. Mag. vol. 24, p.
479, 1887, Electrical Papers, vol. II., p. 355.}
Thus, if a, d be two complex quantities
representing two conductors in series, the corresponding quantity for the
combination is Again, if a, d represent
two conductors in parallel, the reciprocal of the resultant is given by
addition of the reciprocals of a, d. For, if the currents be r, x
corresponding to a difference of potentials V at the common terminals,

ax - dic so that a: + a;' = F(1 ja + 1  a ).

In the application to
Wheatstone’s combination of the general theory of forced vibrations, we mil
limit the impressed forces to the battery and the telephone branches. If x, y
be the currents in these branches, X, Y the corresponding electro-motive
forces, we have, § 107, linear relations between x, y, and A, F, which may be
written A = Ax -f fiy I Y Bx Cy) (3).  the coefficient of y in the first
equation being identical with that of X in the second equation, by the
reciprocal property- The three constants A,ByC are in general complex
quantities, functions of p.  The reciprocal relation may be interpreted as
follows.  Fs= 0, Bx -H Cy = 0, and

 --F=ac If In like manner, if we had supposed X = 0, we should have found

shewing that the ratio of the current in one member to the electromotive force
operative in the other is independent of the way in which the parts are
assigned to the two members.

We have now to determine the constants A, JB. C in terms of
the electrical properties of the system. If y be maintained zero by a suitable
force Y, the relation between x and A' is A' = A therefore denotes the
(generalized) resistance to any electromotive force in the battery member, when
the telephone member iw open. This resistance is made up of j the resistance in
the battery member, and of that of the conductors a -h r, h + f , combined in
parallel. Thus

A-r-   ' j ' 1 4- 0 -h c + a

In like manner ( + 6)(G + rf) n 6 -4 c 4 To determine B let us consider the
force which must act in order that the current through it may be
zero, in spite of the operation of X. We have Y-Br. The total current j flows
partly along the branch la + c), and partly along (6-l-rf).  The current
through (a + c) is

r (a4c)

l ( +c)+ l (t-rrf) H-6 + c + rf

and that through (6 + ci) is

b '4 0 4 d

The difference of potentials at the terminals of e, supposed to he interrupted,
is thus

0 (i + d) 37 - d(a4-c)g a+ft+c+d and accordingly B = - ‘

By (6), (7), (10) the relationship of X, Y to x, y is completely determined.

The problem of the bridge requires the determination of the current y as
proportional to X, when y=0, that is when no electro-motive force acts in the
bridge itself ; and the solution is given at once by the introduction into (4)
of the values ( A, B, C from (6), (7), (10).

If there be an approximate  balance, the expression simplified For (6c -od) is
then small, and B may be neglected relatively to 29-1 AC in the denominator of
(4). Thus, as a sufficient approximation in this case, we may write X~ AC (10)
(C)x(7) (11).  The following interpretation of the process leads very simply to
the approximate form (11), and is independent of the general theory. Let us
first inquire what electro-motive force is necessary in the telephone member to
stop any current through it. If such a force act, the conditions are,
externally, the same as if the member were open ; and the current x in the
battery member due to a force equal to X in that member is XjA, where .4 is
written for brevity ajs representing the right-hand member of 1(6). The
difference of potentials at the terminals of e, still supposed to be open, is
found at once when x is known. It is given by

c X (S)-d X (9 = Ba, where B is defined by (10). In terms of X the difference
of potentials is thus BXjA. If e be now closed, the same fraction expresses the
force necessary in e in order to prevent the generation of a current in that
member.  The case with which we have to deal is when X acts in f and there is
no force in e. We are at liberty, however, to suppose that two opposite
forces, each of magnitude BXjA, act in e. One of these, as we have seen, acting
in conjunction with in , gives no current in c ; so that, since electro-motive
forces act independently of one another, the actual current in e, closed
without internal electro-motive force, is simply that due to the other
component.  The question is thus reduced to the determination of the current in
e due to a given force in that member.

So far the argument is rigorous ; but we will now suppose that we have to deal
with an approximate balance. In this case a force in e gives rise to very
little current in , and in calculating the current in e, we may suppose  to be
broken. The total resistance to the force in e is then given simply by C of
equation (7), and the approximate value for y is derived by dividing BX A by C,
as we found in (11).

A continued application of the foregoing process gives y X in the form of an
infinite geometric series : y - Ail  4


This is the rigorous solution already found in (4); but the first term of the
series suffices for practical purposes.

The form of (11) enables us at once to compare the effects of increments of
resistance and of inductance in disturbing a balance.  For let ad = he, and
then change rf to d -f d where d = rf + id.'.  The value of y A' is
proportional to d and the amplitude of the vibratory current in the bridge is
proportional to
mod d that is to
Thus d are equally efficacious when numerically equal
\footnote{On the Bridge Method in its Application to Periodic Electric Current.
Proc.  Roy. Soc., vol. 40, p. 203, 1891}.
In most cases where a telephone is employed the balance is
more sensitive to changes of inductance than to changes of resistance.

In the use of the Wheatstone balance for purposes of measurement, it is best to
make a equal to c. The equality of h and d can then be tested by interchange of
a and c, independently of the exactitude of the equality of these quantities.
Another advantage lies in the fact that balance is independent of mutual
inductance between a and c or between b and d.

\section{235 q Hughes' Arrangement}
235 q. In the formulae of § 235 p it has been assumed that there is no mutual
induction between the various members of the combination. The more general
theory has been considered very fully by Heaviside
\footnote{On the Self-Induction of Wires, Part VI.; Phil. Mag., Feb. 1887;
Electrical Papers, 1892, vol. II., p. 281.}, but to enter upon it would
lead us too far.  It may be well, however, to sketch the theory of the
arrangement adopted by Hughes, which possesses certain advantages in dealing
with the electrical properties of wires in short lengths \footnote{Journal, Tel.
Eng., vol. XV. (1886) p. 1 ; Proc. Roy. Soc., vol. XL. (1886) p.  451.}

FIG Figure figure Fig. 53 b.

The apparatus consists of a Wheatstone’s quadrilateral. Fig. 53 b, with a
telephone in the bridge, one of the sides of the quadrilateral being the wire
or coil under examination (P), and the other three being the parts into which a
single German-silver wire is divided by two sliding contacts. If the battery
branch (B) be closed, and a suitable interrupter be introduced into the
telephone branch (2), balance may be obtained by shifting the contacts. Provided
that the interrupter introduces no electro-motive force of its own  the balance
indicates the proportionality
of the four resistances. If P be the unknown resistance of the conductor under
test, Q, R the resistances of the adjacent parts of the divided wire, S that of
the opposite part (between the sliding contacts), then, by the ordinary rule.
PS = QR while Q, R, S are subject to the relation
\[Q + R + S= W,\]
W being a constant. If now the interrupter be transferred from the telephone
to the battery-branch, the balance is usually disturbed on account of
induction, and cannot be restored by any mere shifting of the contacts. In
order to compensate the induction, another influence of the same kind must'
be introduced. It is here that the peculiarity of the apparatus lies. A coil
(not shewn in the figure) is inserted in the battery and another in the
telephone-branch which act inductively upon one another, and are so mounted
that the effect may be readily varied. The two coils may be concentric and
relatively movable about the common diameter. In this case the action vanishes
when the planes are perpendicular. If one coil be very much smaller than the
other, the coefficient of mutual induction M is proportional to the cosine of
the angle between the planes. By means of the two adjustments, the sliding of
the contacts and the rotation of the coil, it is usually possible to obtain a
fair silence.

Hughes interpreted his observations on the basis of an assumption that the
inductance of P was represented by A, irrespective of resistance, and that the
resistance to variable currents could (as in the case of steady currents) be
equated to \(QR/S\).  But the matter is not quite so simple. The true formulae are,
however, readily obtained for the case where the only sensible induction among
the sides of the quadrilateral is the inductance L of the conductor P.

Since there is no current through the bridge, there must be the same current
(x) in P and in one of the adjacent sides (say) P, and for a like reason the
same current y in Q and S. The difference of potentials at time t between the
junction of P and R and the junction of Q and S may be expressed by each of the
three following equated quantities :  A condition not always satisfied in
practice with Hughes arrangement.

Introducing the assumption that all the quantities vary harmonically with
frequency \(p 2 \pi\), and eliminating the ratio y : y,
\todo{inline}[smudged]
we find as the conditions required for silence in the telephone
\begin{align}
QR-SP & = p^2 ML \nonumber \\
M(P Q R S) & = S L  \nonumber
\end{align}
(2).  It will be seen that the ordinary resistance balance (SP = QR)
is departed from. The change here considered is peculiar the apparatus
and, so far as its influence is concerned, it does not indicate a real
alteration of resistance in the wire. Moreover since p is involved, the
disturbance depends upon the rapidity of vibration, so that in the case of
ordinary mixed sounds silence can be attained only approximately. Again, from
the second, equation we see that M is mt in general a correct measure of the
value of L \footnote{Discussion on Prof. Hughes’ Address. Journ Tel, Eng, vol,
XV., p. 54  Feb., 1886.}

If, however, P be known, the application of (2) presents no difficulty. In many
cases we may be sure beforehand that P, viz. the effective resistance of the
conductor, or combination of conductors, to the variable currents, is the same
as if the currents were steady, and then P may be regarded as known. But
then' are other cases, - some of them will be alluded to below - in which this
assumption cannot be made; and it is impossible to determine the unknown
quantities L and P from (2) alone. W may then fall back upon (1), By means of
the two equations P and L can always be found in terms of the other quantities.
But among these is included the frequency of vibration ; so that
the method is practically applicable only when the interrupter is
such as to give an absolute periodicity. A scraping contact,
otherwise very convenient, is thus excluded; and this is
undoubtedly an objection to the method.

If the member P be without inductance, but be interrupted by a leyden of
capacity 7, the same formula may be employed, with substitution of - for L.
Equation (1) then gives a measure of G which is independent of the frequency.

235 r. The success of experiments with this kind of apparatus depends very
largely upon the action of the interrupter by which the currents are rendered
variable. When periodicity is not
necessary, a scraping contact, actuated by a clock or by a small motor, answers
very well; but it is advisable, following Lodge and Hughes, so to arrange
matters that the current is suspended altogether at short intervals. The faint
scraping sound, heard in the neighbourhood of a balance, is more certainly
identified when thus rendered intermittent.

But for many of the most interesting experiments a scraping contact is
unsuitable. When the inductance and resistance under observation are rapidly
varying functions of the frequency, it is evident that no sharp results are
possible without an interrupter giving a perfectly regular electrical
vibration. With proper appliances an absolute silence, or at least one
disturbed only a slight sensation of the octave of the principal tone, can be
an ived at under circumstances where a scraping contact would admit of no
approach to a balance at all.

Tuning-forks, driven electromagnetically with liquid or solid contacts (§ 64),
answer well so long as the frequency required does not exceed (say) 300 per
second ; but for experiments with the telephone w e desire frequencies of from
500 to 2000 per second.  Good results may be obtained with harmonium reed
interrupters, the vibrating tongue making contact once during each period with
a stop, which can be adjusted exactly to the required position by means of a
screw .

But perhaps the best interrupter for use with the telephone is obtained by
taking advantage of the instability of a jet of fluid.  If the diameter and the
speed be chosen suitably, the jet may be caused to resolve itself into drops
under the action of a tuning fork in a perfectly regular manner, one drop
corresponding to each complete vibration of the fork. Each drop, as it passes,
may be made to complete an electric circuit by squeezing itself between the
extremities of two fine platinum wires. If the electro-motive force of the
battery be pretty high, and if the jet be salted to improve its conductivity,
sufficient current passes, especially if the aid of a small step-down
transformer be invoked.  Finally the apparatus is made self-acting by bringing
the fork under the influence of an electro-magnet, itself traversed by the same
intermittent current. Such an apparatus may be made to work with frequencies up
to 2000 per second, and it possesses many advantages, among which maybe
mentioned almost absolute
constancy of pitch, and the avoidance of loud aerial disturbance.  The
principles upon which the action of this interrupter depends will be further
considered in a subsequent chapter.

Phil. Mag,, vol. 22, p. 472, 1886

235 s. Scarcely less important than the interrupter are the arrangements for
measuring induction, whether mutual induction, as required in § 235 q, or
self-induction. Inductometers, as Heaviside calls them, may be conveniently
constructed upon the pattern of Hughes. A small coil is mounted so that one
diameter coincides with a diameter of a larger coil, and is movable about that
diameter. The mutual induction M between the two circuits depends upon the
position given to the smaller coil, which is read by a pointer attached to it,
and moving over a graduated circle. If the smaller coil were supposed to be
infinitely small, the value of iV, as has already been stated, would be
proportional to the sine of the displacement from the zero position (M = 0).
But an approximation to this state of things is not desirable. If the mean
radius of the small coil be increased until it amounts to 55 of that of the
larger, not only is the efficiency much enhanced, but the scale of Jl is
brought to approximate coincidence, over almost the whole practical range, with
the scale of degrees \footnote{Phil Mag., vol. 22. p. 498. 1886.}

The absolute value of each degree may be arrived at in
various ways, perhaps most simply by adjusting the mutual induction of the
instrument to balance a standard of mutual induction.

For experiments upon the plan of § 235 q the one coil is included in the
telephone and the other in the battery branch, but when the object is to secure
a variable and measurable inductance, the two coils are connected in series.
The inductance of the combination is then \(L+2M+N\) of which the first and
third terms are independent of the relative position of the coils.

\section{~~235 t}
235 t. Good results by the method of § 235 q have been obtained by Weber
\footnote{Electrical Review, April 9, July 9, 1886}, and by
the author\footnote{Phil, Mag., loc. cit.}
using a reed interrupter of frequency 1050 per second i but the fact
that inductance and resistance are mixed up in the measurements is a decided
drawback, if it be only because the readings require for their interpretation
calculations not readily made upon the spot.

The more obvious arrangement is one in which both the induction and the
resistance of the branch containing the subject under examination are in every
case brought up to the given totals necessary for a balance. To carry this out
conveniently we require to be able to add inductance without altering
resistance, and resistance without altering inductance, and both in a
measurable degree. The first demand is easily met. If we include in the
circuit the two coils of an inductometer, connected in series, the inductance
of the whole can be varied in a known manner by rotating the smaller coil. On
the other hand the introduction, or removal, of resistance without alteration
of inductance cannot well be carried out with rigour. But in most cases the
object can be sufficiently attained with the aid of a resistance-slide of thin
German-silver wire which may be in the form of a nearly close loop.

In the Wheatstone s quadrilateral, as arranged for these experiments, the
adjacent sides It, S may be made of similar wires of German silver of equal
resistance ohm). If doubled they give rise to little induction, but the
accuracy of the method is independent of this -circumstance. The side P
includes the conductor, or combination of conductors, under examination, an
inductometer, and the resistance-slide. The other side. Q, must possess
resistance and inductance greater than any of the conductors to be compared,
but need not be susceptible of ready and measurable variations. In order to
avoid mutual induction between the branches, P and Q should be placed at same
distance away; being connected with the rest of the apparatus by leads of
doubled wire.

It will be evident that when the interrupter acts in the battery branch,
balance can be obtained at the telephone in the bridge only under the
conditions that both the inductance and the resistance in P are equal in the
aggregate to the corresponding quantities in Q. Hence when one conductor is
substituted for another in P, the alterations demanded at the inductometer and
in the slide give respectively the changes of inductance and of resistance. In
this arrangement inductance and resistance are well separated, so that the
results can be interpreted without calculation; but the movable contacts of the
slide appear to introduce uncertainty into the determination of resistance.

In order to get rid of the objectionable movable contacts some sacrifice of
theoretical simplicity seems unavoidable. We can no longer keep the total
resistances P and Q constant; but by reverting to the arrangement adopted in a
well-known form of Wheatstone s bridge, we cause the resistances taken from P
to be added to Q, and vice versa. The transferable resistance is that of a
straight wire of German-silver, with which one telephone terminal makes contact
at a point whose position is read off on a
divided scale. Any uncertainty in the resistance of this contact does not
influence the measurements.

The diagram Fig. ( 53 c) shows the connection of the parts. One
of the telephone terminals T goes to the junction of the only
resistances li and S, the other to a point upon the divided wire.
The branch P includes one inductometer (with coils connected in
series), the subject of examination, and part of the divided wire.
The branch Q includes a second inductometer (replaceable by a
simple coil possessing suitable inductance), a rheostat, or any
resistance roughly adjustable from time to lime, and the remainder pf the
divided wire. The battery branch B, in which may
also be included the interrupter, has its terminals connected, one
to the junction of P and li, the other to the junction of Q and H.
When it is desired to use steady currents, the telephone can of
course be replaced by a galvanometer.

In this arrangement, as in the other, balance requires that the
branches P and Q be similar in respect both of inductance and of
resistance. The changes in inductance dm to a shift in the
movable contact may usually be disregarded, and thus any alteration in the
subject (included in P) is measured by the rotation
necessitated at the inductometer. As for the resistance, it is
evident that (K and S being equal) the value for any additional
conductor interposed in P is measured by twice the displacement
of the sliding contact necessary to regain the balance.

Experimental details of the application of this method to the
measurement of various combinations will be found in the paper'
from which the above sketch is derived. Among these may be
mentioned the verification of Maxwell’s formulae, (8), (9) § 285 k,
as to the influence of a neighbouring circuit, especially in the
extreme case where the equivalent inductance is almost destroyed,
and of the formula (10) § 235 m relating to the behaviour of an
electro-magnet shunted by a relatively high simple resistance.
But the most interesting in many respects is the application to
the phenomena presently to be considered, where the conductors
in question are no longer approximately linear but must be
regarded as solid masses in which the currents are distributed in
a manner that needs to be determined by general electrical theory.

As has already been remarked more than once, a ley den may
always be supposed to be included in the circuit, the stiffness
, thereof having the effect of a negative inductance. If there be no
hysteresis in the action of the leyden, the whole effect is thus
represented ; but when the dielectric employed is solid, it appears
that dissipative loss cannot be avoided. The latter effect manifests
itself as an augmentation of apparent resistance, indistinguishable,
unless the frequency be varied, from the ordinary resistance of the
leads. A similar treatment may be applied to an electrolytic cell,
the stiffness and resistance being presumably both functions of the
frequency.

\section{~~235 u}
235 u. It was proved by Maxwell that a perfectly conducting sheet, forming a
closed or an infinite surface, acts as a magnetic screen, no magnetic actions
which may take place on one side of the sheet producing any magnetic effect on
the other side.  In practice we cannot use a sheet of perfect conductivity; but
the above described state of things may be approximated to in the case of
periodic magnetic changes, if the time-constants of the sheet circuits be large
in comparison with the periods of the changes.

 Phil, loc. eit.  Electricity and Magnetism, 1673, § 655.

The experiment is made by connecting up into a primary
circuit a battery, a microphone clock, and a coil of insulated wire.
The secondary circuit includes a parallel coil and a telephone.
Under these circumstances the hissing sound is heard almost as
well as if the telephone were inserted in the primary circuit
itself. But if a large and stout plate of copper bo interposed
between the two coils, the sound is greatly enfeebled. By a proper
choice of battery and of the distance between the coils, it is not
difficult so to adjust the strength that the sound is conspicuous in
the one case and inaudible in the other

One of the simplest applications of Maxwell's principle is to
the case of a long cylindrical shell placed within a coaxial magnetizing helix.
The condition of minimum energy requires that such
currents be developed in the shell as shall neutralize at internal
points the action of the coil. Thus, if the conductivity of the
shell be sufficiently high, the interior space is screened from the
magnetizing force of periodic currents flowing in the outer helix,
and conducting circuits situated within the shell must be devoid
of induced currents. An obvious deduction is that the currents
induced in a solid conducting core will more and more confined
to the neighbourhood of the surface as the frequency of electrical
vibration is increased.
The point at which the concentration of current towards the
surface becomes important depends upon the relative values v of the
imposed vibration-period aid the principal time-constant of the
core circuit. If p be the specific resistance of the material, its
magnetic permeability, the radius of the cylinder, the expression
for the induction (c) parallel to the axis, during the progress of the
subsidence of free currents in a normal mode, is

c =  Ji,(kr) (1)
where
47rX i
(2),
and ka is determined by the condition that

J (ka)-0 (3).

The roots of (3) are, § 206, 2-404, 5-520, 8 654, 11-792, etc.,
so that for the principal mode of greatest persistence 2-404 r a) (4),
 2-404yp

Acoustical Observations, Phil. Mag, voL 13, p. 344, 1882.
For copper in C.G.S., measure p = 1642, 1, and thus 800 nearly.
In the case of iron we may take as approximate values, p, = 100,
p = 10‘. Thus for an iron wire of diameter (2a) equal to 33 cm.,
the value of t is about of a second, and is therefore comparable
with the periods concerned in telephonic experiments.
Regarded from an analytical point of view- the theory of forced
vibrations in a conducting, core is equally simple, and was worked
out almost simultaneously by Lamb, Oberbeck and Heaviside .
In this case we are to regard  as given, equal (say) to ip. where
pfiir is the frequency. If be the imposed magnetizing force,
the solution is

c J ihr) JuQca) (6),
the value of k being given by (2).

When the period in the field is long in comparison with the
time of decay of free currents, we have Jro(ir) = 1, nearly, so that
c is approximately constant and throughout the section of
the cylinder. But, in the opposite extreme, when the oscillations
in the intensity of the field are rapid in comparison with the decay
of free currents, the induced currents extend only to a small depth
beneath the surface of the cylinder, the inner strata (so to speak)
being almost completely sheltered from electromotive force by the
outer ones. Writing = (1 - f), aji here
we have, when qr is large,

t o(T) = const. X,

Vr

approximately, and thence

c = ft . - (a r) . eS' + (r-at j

This indicates that the electrical disturbance in the cylinder
consists in a series of waves propagated inwards with rapidly
diminishing amplitude'.

the Duration of Free Electric Currents in an Infinite Conducting Cylinder,
Brit. Jmoc. Report for 1862, p. 446.

Proe. Math. Soc., vol. xv., p. 189, Jan. 1884.

Jfied. Ann,t vd. xzi., p. 672, Ap. 1884.
Electrician May, 1884. Electrical Papers vol. II., p. 353.

For experimental purposes what we most require to know is
the reaction of the core currents upon the helix, in which alone
we can directly measure electrical effects. This problem is fully
treated by Heaviside, but we must confine ourselves here to a
mere statement of results. These are most conveniently expressed
by the changes of effective inductance L and resistance R due to
the core. If m be the number of turns per unit length in the
magnetizing helix, and if SZ, hR be the apparent alterations of L
and R due to the introduction of the core, also reckoned per unit
length, we have

hL =4w' 7ra2( iP- 1))

hR = 4 r7r‘a'‘‘ i )Q 1

where P and Q are defined by

= (8).

the function f being of the form
and the argument x being

(10).

If the material composing the cor s be non-conducting, x = 0, and
therefore

p=i, e=o.
Accordingly = 4m-7r a= (ft - 1), SR = 0 (11).

These values apply also, whatever be the conductivity of the
core, if the frequency be sufficiently low.
At the other extreme, when  we require the ultimate
form of from the value of Jo given in (10) § 200, or other
wise, it may be shewn that in the limit
so that
=

p o=
(12),

(13).
The introduction of these values into (7) shews that in the
limit, when the frequency is exceedingly high,

8i - - 4 mi r a, SB  0 (14),
' Lamb, loc. eit., where is also discussed the problem of the current induced
by the sudden cessation of a previously constant field.

 loc. cit.
as might also have been inferred from the consideration that the
induced currents are then confined to the surface of the core.

An example of the application of these formulae to an intermediate case and a
comparison with experiment will be found in the paper already referred to.

\section{~~235 v}
235 v. The application of Maxwell’s principle to the case of
a wire, in which a longitudinal electric current is induced, is less
obvious; and Heaviside appears to have been the first to state
distinctly that the current is to be regarded as propagated inwards
from the exterior. The relation between the electromotive force
E and the total current C had, however, been given many ears
earlier by Maxwell in the form of a series. His result is equivalent to
E

RC
(1),
in which R denotes the whole resistance of the length I to
steady currents, fi the permeability, and p 27r the frequency. The
function  is that defined by (9) § 235 te, and is a constant
dependent upon the situation of the return currents
The most convenient form of the results is that which we have
already several times employed. If we write
E R'C ipL'C
(2),
in which K and L' are real, these quantities will represent the
effective resistance and inductance of the wire. When the argument in (1) is
small, that is when the frequency is relatively low,
we thus obtain
(3),

I'll - X t - + ...) (4)‘.


1 Phil Mag., vol. 22, p. 493, 1886.

- Electrician, Jan., 1885 ; Electrical Papers, vol. i., p. 440.

 Phil, Trans., 1865 ; Electricity and Magnetism, vol. II., § 690.

 The simplest case arises when the dielectric, which bounds the cylindrical
wire of radius a, is enclosed within a second conducting mass extending outwards
to infinity and bounded internally at a cylindrical surface r=b. We then have
if = 2 log (hja). See J. J. Thomson, loc. cit., § 272.

3 Phil. Mag., vol. 21, p. 887, 1886.

It is singular that Maxwell (Zoc. cit.) seems lo have regarded his solution as
conveying a correction to the self-induction only of the win.

.When p ia very small, these equations give, as was to be
expected,

 = -R, L (a).

If we include the next terms, we recognise that, in accordance
with the general rule, L begins to diminish and R' to increase.

When p is very great, we have to make use of the limiting
form of
\( \phi' /  \phi\).
As in § 235 u,
\[ \phi' /  \phi = (1 + i) \sqrt{ p l \mu /(2 R) }, \]
(6) and thus ultimately
\[ R'  = \sqrt{ p l \mu /(2 R) }, \]
(7),
\todo[inline]{smudged}
\[ L'/l = A  + \sqrt{\mu R 2 pl }\]

(8), the first of which increases without limit with p, while the second
tends to the finite limit A, corresponding to the total exclusion of
current from the interior of the wire.

Experiments \footnote{Phil. Mag. vol. 22, p. 488, 1886.}
upon an iron wire about 18 metres long and
millimetres in diameter led to the conclusion that the resistance
to variable currents of frequency 1050 was such that l O.

A calculation based upon (1) shewed that this result is in harmony
with theory, if x = 99'5. Such is about the value indicated by
other telephonic experiments.

\section{~~235 w Limiting forms}
235 w. The theory of electric currents in such wires as are
commonly employed in laboratory experiments is simple, mainly in
consequence of the subordination of electrostatic capacity. When
this element can be neglected, the current is necessarily the same
at all points along the length of the were, so that whatever enters
a wire at the sending end leaves it unimpaired at the receiving
end. In this case the whole electrical character of the wire can
be expressed by two quantities, its resistance R and inductance Z,
and these may usually be treated as constants, independent of the
frequency. The relation of the current to the electromotive force
under such circumstances has already been discussed (7) § 23-5 j.
When we have occasion to consider only the amplitude of the
current, irrespective of phase, we may regard it as determined
by a quantity which is called by Heaviside the
impedance. Thus in circuits devoid of capacity the impedance is
always increased by the existence of L,


Circuits employed for practical telephony may often be regarded as coming
under the above description, especially when
the wires are suspended and are of but moderate length. But
there are other cases in which electrostatic capacity is the dominating
feature. The theory of electric cables was established
many years ago by Lord Kelvin
\footnote{Proc. Roy. Soc., 1855; Mathematical and Physical Papers vol. II. p. 61.}
for telegraphic purposes. If S
be the capacity and i the resistance of the cable, reckoned per
unit length, V and C the potential and the current at the point we have

SdVdt = -dCd::, I(C=-dV;ds (1),

whence W dC dt = d-C:dz y (2),

the well known equation for the conduction of heat discussed by
Fourier. On the assumption that C is proportional to it
reduces to

d: Cidz ‘= W(ipRS).(l (3);

so that the solution for waves propagate in the positive direction is

(4).

The distance in traversing which the current is attenuated in the
ratio of e to 1 is thus

z s (2 pi RS) (5).

A very slight consideration of the magnitudes involved is
sufficient to give an idea of the difficulty of telephoning through a
long cable. If, for example, the frequency (p 27r) be that of a
note rather more than an octave above middle c and the cable be
such as are used to cross the Atlantic, we have in e Gs measure

v = (0, iRS)- = 2x 10,

and accordingly from (5)

z = 3 X 10 cm. = 20 miles approximately.

A distance of 20 miles would thus reduce the intensity of
sound, measured by the square of the amplitude, to about a
tenth, an operation which could not be repeated often without
rendering it inaudible. With such a cable the practical limit
would not be likely to exceed fifty miles, more especially as
the easy intelligibility of speech requires the presence of tones
still higher than is supposed in the above numerical example.
\footnote{On Telephoning through a Cable. Brit. Assoc. Report for 1884, p. 632.}

Proc. Roy. Soc., 1855; Mathematical and Physical Papers vol. II. p. 61.


Heaviside's theory.
\section{~~235 x Heaviside's Theory}
235 x. In the above theory the instillation is supposed to be
perfect and the inductance to be negligible. It is probable that
these conditions are sufficiently satisfied in the case of a cable,
but in other telephonic lines the inductance is a feature of great
importance. The problem has treated with full generality by Heaviside, but a
slight sketch if his investigation is all that our limits permit.

If I, S, Z, K be the resistance, capacity or permittance, inductance, and
leakage-conductance respectively per unit of length,
V and C the potential-difference and current at distance c, the
equations, analogous to (1) are
\[ KV + S = C, RC + L = -V\]
(1) Thus, if the currents are harmonic, proportional to e,
\[ C = (iR + ipL ) (K + ipS) C\]
(2), with a similar equation for V.

It might perhaps have been expected that a finite leakage \(K\)
would always act as a complication; but Heaviside
\footnote{Electrician, June 17. 1887. Electrical Papers, vol. II. pp. 125, 300}.
has shewn that it may be so adjusted so to simplify the matter. This case,
which is remarkable in itself and also serves to throw light upon
the general question, arises when \(RL=KS\). We will write

LSv = 1, H;L = KjS (2),

where v is a velocity of the order of the velocity of light. The
equation fur V is then by ( 1 )

f d' V I dz' kihdi + qfV ( 3 ) ;

or if we take U so that

F=e l7 (4),

v d U;dz = d Uidr (o),

the well-known equation of undisturbed wave propagation § 144.
Thus, if the wave be positive, or travel in the direction of
increasing z, we shall have, if z be the state of V initially,

V, = fAz - vt), C =  (6).

If Fo, Cn be a negative wave, travelling the other way,

F = + C, VJLc (7).



Thus, any initial state being the sum of Vy and V. to make F,
and of Cy and C. to make C, the decomposition of an arbitrarily
given initial state of V and C into the waves is effected by

F, = i ), F, = i ( vLC) (8).

We have now merely to move Vy bodily to the right at speed
V, and Fj, bodily to the left at speed v, and attenuate them to the
extent to obtain the state at time t later, provided no changes
of condition have occurred. The solution is therefore true for all
future time in an infinitely long circuit. But when the end of a
circuit is reached, a reflected wave usually results, which ' usl be
added on to obtain the real result.”
As in § 144, the precise character of the reflection depends
upon the terminal conditions. One case is uniquely simple.
Let there be a resistance inserted of amount vL. It introduces
the condition V = vLC if at say B, the positive end of the circuit,
and V -vLC if at the negative end, or beginning. These are
the characteristics of a positive and of a negative wave respectively ; it
follows that any disturbance arriving at the resistance is
at once absorbed. Thus, if the circuit be given in any state
whatever, without impressed force, it is wholly cleared of electrification and
current in the time Ijv at the most, if I be the length
of the circuit, by the complete absorption of the two waves into
which the initial state may be decomposed.”

 Blit let the resistance be of amount Ry at say B ; and let F,
and Fa be corresponding elements in the incident and reflected
waves. Since we have
Vy -vLGy, Fa = -vLa,
F,+ Fa = iJ,(C, + C,)...(9),
we have the reflected w ave given by
V.,R jvL
(10).
If Ry be greater than the critical resistance of complete absorption, the
current is negatived by reflection, whilst the electrification does not change
sign. If it be less, the electrification is
negatived, whilst the current does not reverse.”

Two cases are specially notable. They are those in which
there is no absorption of energy. If jBi = 0, meaning a short
circuit, the reflected wave of F is a perverted and inverted copy of

the incident. But if iJ = x, representing insulation, it is C that
is inverted and perverted\footnote{Heaviside, Collected Works, vol. II. p. 312.}

The cases last mentioned are evidently analogous to the reflection of L
sonorous aerial wave travelling in a pipe. If the end of
the pipe be closed, the reflection is of one character, and if it be
open of another character. In both cases the whole energy is
reflected, § 257.
\todo[inline]{ section 257 of what?}
The waves reflected at the two ends of an electric
circuit complicate the general solution, especially when the simplifying
condition (2) does not hold. But in many cases of practical interest they inay
be omitted without much loss of
accuracy. One passage over a long line usually introduces considerable
attenuation, and then the effect of the reflected wave,
which must traverse the line three times in all, becomes insignificant.

In proceeding to the general solution of (2) for a positive
wave, we will introduce, after Heaviside, the following abbreviations,

i LS = 1, R;Lp = , KSp g (11).

In terms of these quantities (2) may be written

((Cldz = (P + iQYC (12).

where

P or ( = i (p vf i(l (1 4. ± (fg  1)) ... (13).

Thus, if P and Q be taken positively, the solution for a wave
travelling in the positive direction is

C = cos ( pt - Qz) (14),

the current at the origin being C cos

The cable formula, § 235 w, is the particular case arrived at by
supposing in (13) = 00 , y = 0, which then reduces to

= (15).

Again, the special case of equation (3) is derivable by putting
 =  = qjp. The result is

P = jI% Q pIv (16).

If the insulation be perfect, g = 0, and (13) becomes

P or 0 ((1 + ) T 1) (17).

In certain examples of long copper lines of high conductivity,
f in y be regarded as small so far as telephonic frequencies are
concerned. Equation (17) then gives

P -pfl v = Rj vL, Q=:pjv (18).

For a further discussion of the various cases that may arise
the reader must be referred to the writings of Heaviside already
cited. The object is to secure, as far as may be, the propagation
of waves without alteration of type. And here it is desirable to
distinguish between simple attenuation and distortion. If, as in
(16) and (18), P is independent of p, the amplitudes of all components are
reduced in the same ratio, and thus a compl wave
travels without distortion. The cable formula (15) is an example
of the opposite state of things, where waves of high frequency are
attenuated out of proportion to waves of low frequency. It appears
from Heaviside’s calculations that the distortion is lessened by
even a moderate inductance.

The effectiveness of the line requires that neither the attenuation nor the
distortion exceed certain limits, which however it is
held to lay down precisely. A considerable amount of distortion
is consistent with the intelligibility of speech, much that is
imperfectly rendered being supplied by the imagination of the hearer.

\section{~~235 y}
235 y. It remains to consider the transmitting and receiving
appliances. In the early days of telephony, as rendered practical
by Graham Bell, similar instruments were employed for both
purposes. Bell’s telephone consists of a bar magnet, or battery
of bar magnets, provided at one end with a short pole-piece
which serves as the core of a coil of fine insulated wire. In close
proximity to the outer end of the pole-piece is placed a circular
disc of thin iron, held at the circumference. Under, the influence
of the permanent magnet the disc is magnetized radially, the
polarity at the centre being of course -opposite to that of the
neighbouring end of the steel magnet.

The operation of the instrument as a transmitter is readily
traced. When sonorous waves impinge upon the disc, it responds
with a symmetrical transverse vibration by which its distance
from the pole-piece is alternately increased and diminished.
When the interval is diminished, more induction passes through
the pole-piece, and a corresponding electro-motive force acts in
the enveloping coil. The periodic movement of the disc thus
gives rise to a periodic current in any circuit connect e  with the
telephone coil.

The electro-motive force is in the first instance proportional
to the permanent magnetism to which it is due; and this law
would continue to hold, were the behaviour of the pole-piece and
of the disc conformable to that of the ‘'soft iron ’’ of approximate
theory. But as the magnetism rises, and the state of saturation
is more nearly approached, the response to periodic changes of
force becomes feebler, and thus the efficiency falls below that
indicated by the law of proportionality. If we could imagine the
state of saturation in the pole-piece to le actually attained, the
induction through the coil would become almost incapable of
variation, being reduced to such as might occur were the iron
removed. There is thus a point, dependent upon the properties
of magnetic matter, beyond which it is pernicious to raising the
amount of the permanent magnetism ; and this point marks the
maximum efficiency of the transmitter. It is probable that the
most favourable condition is not fully reached in instruments
provided with steel magnets; but tin considerations above
advanced may serve to explain why an electro-magnet is not substituted.

The action of the receiving instrument may be explained on
the same principles. The periodic current in the coil automatically
opposes an J cooperates with the permanent magnet, and thus the
iron disc is subjected to a periodic force acting at it ' centre.
The vibrations are thence communicated to the air, and so roach
the ear of the observer. As in the case of the transmitter, the
efficiency attains a maximum when the magnetism of the pole-piece is still far
short of saturation.

The explanation of the receiver in terms of magnetic forces
pulling at the disc is sometimes regarded as inadequate or even as
altogether wide of the mark, the sound being attributed to molecular
disturbances in the pole-piece and disc. There is indeed
every reason to suppose that molecular movements accompany
the change of magnetic state, but the question is how do these
movements influence the ear. It would appear that they can do
so only by causing a transverse motion of the surface of the disc,
a motion from which nodal subdivisions are not excluded.

In support of the  push and pull theory  it may be useful tu
cite in experiment tried upon a bipolar telephone. In this
instrument each end of a horse-shoe magnet is provided with a
pole-piece and coil, and the two pole-pieces are brought into
proximity with the disc at places symmetrically situated with
regard to the centre. In the normal use of the instrument the
two coils are permanently connected as in an ordinary horse-shoe
electro-magnet, but for the purposes of the experiment provision
was made whereby one of the coils could be reversed at pleasure
by means of a reversing key. The sensitiveness of the telephone
in the two conditions was tested by including it in the circuit of
a Daniel cell and a scraping contact apparatus, resistance room a
box being added until the sound was but just easily audible.
The resistances employed were such as to dominate the self-induction of the
circuit, and the comparison shewed that the
reversal of the coil from its normal connection lowered the sensitiveness to
current in the ratio of 11 : 1. That the reduction was
not still greater is readily explained by outstanding failures of
symmetry; but on the molecular disturbance theory it is not
evident why there should be any reduction at all.

Dissatisfaction with the ordinary theory of the action of a
receiving telephone may have arisen from the difficulty of understanding how
such very minute motions of the plate could be
audible. This is, however, a question of the sensitiveness of the
ear, which has been proved capable of appreciating an amplitude
of less than 8 x 1O cm.  \footnote{Proc. Roy. Soc. vol. XXVI. p. 248, 1877.}
The subject of the audible minimum
will be further considered in the second volume of this work.


The calculation a priori of the minimum current that should
be audible in the telephone is a matter of considerable difficulty ;
and even the determination by direct experiment has led to
widely discrepant numbers. In some recent experiments by the
author a unipolar Bell telephone of 70 ohms resistance was
employed. The circuit included also a resistance box and an
induction coil of known construction, in which acted an electromotive force
capable of calculation. Up to a frequency of 307
this could be obtained from a revolving magnet of known moment
and situated at a measured distance from the induction coil.  For
the higher frequencies magnetized tuning-forks, vibrating with
measured amplitudes, were substituted. In either case the
resistance of the circuit was increased until the residual sound
was but just easily audible. Care having been taken so to
arrange matters that the self-induction of the circuit was negligible, the
current could then be deduced from the resistance and
the calculated electro- motive force operating in the induction
coil. The following are the results, in which it is to be understood that the
currents recorded might have been halved without
the sounds being altogether lost :
Pitch Source

The effect of a given current depends, of course, upon the
manner in which the telephone is wound. If the same space be
occupied by the copper in the various cases, the current capable of
producing a particular effect is inversely as the square root of the
resistance.

The numbers in the above table giving the results of the
author’s experiments are of the same order of magnitude as
those found by Ferrari
\footnote{Atti della Accad. d. Sci. Di Torino, vol. XIII. p. 1024, 1877.}
whose observations, however, related
to sounds that were not pure tones. But much lower estimates
have been put forward. Thus Tait
\footnote{Edin. Proc. vol.  p. 551, 1878.}
gives \(2 \times 10^{-12}\)  amperes,
and Preece a still lower figure,
\(6 \times 10^{-13}\).  These discrepancies
enormous as they stand, would be still further increased were
the comparison made to refer to the amounts of energy absorbed.

According to the calculations of the author the above tabulated
sensitivity to a periodic current of frequency 256 is about what
might reasonably be expected on the push and pull theory
\footnote{I propose shortly to publish these calculations.}.
At this frequency, which is below those proper to the telephone plate
(§ 22i a), the motion of the plate is governed by elasticity rather
than by inertia, and an equilibrium theory (§ 100) is applicable as
a rough approximation. The greater sensitiveness of the telephone
at frequencies in the neighbourhood of .512 would appear to
depend upon resonance (§ 46). It is doubtful whether the much
higher sensitiveness claimed by Tait and Preece could be reconciled with theory.

It appears to be established that the iron plate of a telephone
may be replaced by one of copper, or even of non-conducting
material, without absolute loss of sound ; but these effects are
probably of a different order of magnitude. In the case of copper
induced currents may confer the necessary magnetic properties.
For a description of the ingenious receiver invented by Edison
and for other information upon telephonic appliances the reader
may consult Preece and Stubbs’ Manual of Telephony,

In existing practice the transmitting instrument depends
 Hon a variable contact. The first carbon transmitter was constructed by Edison
in 1877, but the instruments now in use are
modifications of Hughes microphone
\footnote{Proc. Roy, Soc., vol. XXVII. p. 362, 1870.}.
A battery current is led
into the line through pieces of metal or of carbon in loose juxtaposition,
carbon being almost universally employed in practice.
Under the influence of sonorous vibration the electrical resistance
of the contacts varies, and thus the current in the line is rendered
representative of the sound to be reproduced at the receiving end.

That the resistance of the contact should vary with the
pressure is not surprising. If two clean convex pieces of metal
are forced together, the conductivity between them is represented
by the diameter of the circle of contact (§306). The relation
between the circle of contact and the pressure with which the
masses are forced together has been investigated in detail by
Hertz \footnote{Crelle, Journ, Math. xcn. p. l.'JS, 1882.}.
His conclusion for the case of two equal spheres is that
the cube of the radius of the circle of contact is proportional to
the pressure and to the radii of the spheres. But it has not yet
been shewn that the action of the microphone can be adequately
explained upon this principle.



\chapter{APPENDIX.  On Progressive Waves.} % 476-480

From the Proceedings of the London Mathematical Society Vol IX, p. 21, 1877.

It has often been remarked that, when a group of waves advances
into still water, the velocity of the group is less than that if the individual
waves of which it is composed ; the waves appear to advance
through the group, dying away as they approach its anterior limit.
This phenomenon was, I believe, first explained by Stokes, who regarded the
group as formed by the superposition of two infinite trains
of waves, of equal amplitudes and of nearly equal wave-lengths, advancing in
the same direction. My attention was called to the subject
about two years since by Mr Froude, and the same explanation then
occurred to me independently
\footnote{Another phenomenon, also mentioned to me by Mr Fronde, admits of ft similar
explanation. A steam-launch moving slowly through the water is accompanied by
a peculiar system of diverging waves, of which the most striking feature is the
obliquity of the line containing the greatest elevations of successive waves to the
wave-fronts. This wave pattern may be explained by the superposition of two
(or more) infinite trains of waves, of slightly differing wave-lengths, whose directions
and velocities of propagation are so related in each case that there is no change of
position relatively to the boat. The mode of composition will be best understood by
drawing on paper two sets of parallel and equidistant lines, subject to the above
condition, to represent the crests of the component trains. In the case of two trains
of slightly different wave-lengths, it may be proved that the tangent of the angle
between the line of maxima and the wave-fronts is half the tangent of the angle
between the wave-fronts and the boat’s course.}. In my book on the
\textit{Theory of Sound} (§191),
I have considered the question more generally, and
have shewn that, if V be the velocity of propagation of any kind of
waves whose  wave-length is
\(\lambda\),
and
\(k=2 \pi/\lambda\), then \(U\), the velocity of
a group composed of a great number of weaves, and moving into an undisturbed
part of the medium, is expressed by
\[ U = \frac{d}{dk} ( kV ) \],
(1), or, as we may also write it,

(2) Thus, if \(V \propto \lambda^n \), \( U = (1-n) V\) (3).

In fact, if the two infinite trains be represented by cos k(Vt - x)
and cos K ( V't - x), their resultant is represented by

cos k(Vt-x) + cos k' ( V't - ac),

which is equal to
 2 cos j- 2-‘--2 j- . cos kT' + kV k' + k 2
If k'-kf V'- V be small, we have a train of waves whose amplitude
varies slowly from one point to another between the limits 0 and 2,
forming a series of groups separated from one another by regions comparatively
free from disturbance. The position at time t of the middle
of that group, which was initially at the origin, is given by
(k'r--kF)t (k'-k)x=0,

which shews that the velocity of the group is (K V  kV) -r (K k).
In 'the limit, when the number of waves in each group is indefinitely
great, this result coincides with (1).  The following particular cases are
worth notice, and are here tabulated for convenience of comparison :

FflcX, Reynolds’ disconnected pendulums.

TocX , Deep-water gravity waves.

FxX , U V, Aerial waves, Ac.

FxX-i, II Capillary water waves.

FxX-‘, Z7=2F, Flexural waves.

The capillary water waves are those whose wave-length is so small
that the force of restitution due to capillarity largely exceeds that due
to gravity. Their theory has been given by Thomson (Phil. Mag., Nov. 1871).
The flexural waves, for which Z7=2P, are those corresponding to the
bending of an elastic rod or plate (Theory of Sound. § 191).

In a paper read at the Plymouth meeting of the British Association
(afterwards printed in Nature, Aug. 23, 1877), Prof. Osborne
Reynolds gave a dynamical explanation of the fact that a group of
deep-water waves advances with only half the rapidity of the individual waves.
It appears that the energy propagated across any point,
when a train of waves is passing, is only one-half of the energy necessary to
supply the waves which pass in the same time, so that, if the
train of waves be limited, it is impossible that its front can be propagated
with the full velocity of the weaves, because this would imply the
acquisition of more energy than can in fact lx supplied. Prof. Reynolds
did not contemplate the cases where more energy is propagated than
corresponds to the waves passing in the same time ; but his argument,
applied conversely to the results already given, shews that such cases
must exist. The ratio of the energy propagated to that of the missing
waves is C : V ; thus the energy propagated in the unit time is: T
of that existing in a length T, or U times that existing in the unit
length. Accordingly

Energy propagated in unit time : Energy contained (on an average)
in unit length \( =d (kV) : dk\) by (1).

As an example, T will take the case of small irrotational weaves in
water of finite depth
\(\ell\)
\footnote{Prof. Reynolds considers the trochoidal wave of Rankine and Fronde, which
involves molecular rotation.}.
If z be measured downwards from the surface,
and the elevation (h) of the wave lie denoted by

  - 1 cos (nf - kx) ( 1 ),
in which n - kV the corresponding velocity-potential is
r  - in (j).
This value of satisfies the general differential equation for irrotational motion
v makes the vertical velocity zero when
s = , and - when z -- 0. The velocity of propagation is given by
f F-: fh).
We may now calculate the energy contained in a length as, which is
supposed to include so great a number of waves that fractional parts
may be left out of account.

For the potential energy we have
jz dz dx yp jP dx- gpl . x (7).

For the kinetic energy, (8),
by (1) and (6). If, in accordance with the argument advanced at the
end of this paper, the equality of T, and T be assumed, the value of
the velocity of propagation follows from the present expressions. The
whole energy in the waves occupying a length x is therefore
(for each unit of breadth) F, + T - -x (9),

H denoting the maximum elevation.

We have next to calculate, the energy propagated in time t across a
plane for which x is constant, or, in other words, the work (11) that
must be done in order to sustain the motion of the plane (considered
as a flexible lamina) in the face of the fluid pressures acting upon the
front of it. The variable part of the pressure (Sp), at depth 2:, is
given by

at  1
while for the horizontal velocity
the
= kVU
cos (nt - a:) ;
so that = (10),

on integration. From the value of V in (6) it may be proved that

di =V r- tt ( 4'

and it is thus verified that the value of W for a unit time
d(kV)

dk
X
energy in unit length.

As an example of the direct calculation of U, we may take the case
of waves moving under the joint influence of gravity and cohesion.

It is proved by Thomson that
(11),
where T' is the cohesive tension. Hence
(12).

When k is small, the surface tension is negligible, and then \(U=V/2\)
but when, on the contrary, k is large, as has already been
stated. When TJd = y, i7 = F. This corresponds to the minimum
velocity of propagation investigated by Thomson.


Although the argument from interference groups seems satisfactory,
an independent investigation is desirable of the relation bet i en
energy existing and energy propagated. For some time I was at a loss
for a method applicable to all kinds waves, not seeing in particular
why the comparison of energies should introduce the consideration of
a variation of wave-length. The following investigation, in which the
increment of wave-length is imaginary may perhaps be considered to
meet the want.

Let us suppose that the motion of every part of the medium is
resisted by a force of very small magnitude proportional to the mass
and to the velocity of the part, the effect of which will be that waves
generated at the origin gradually die away as x increases. The motion,
which in the absence of friction would be represented by cos (nt - X-r),
under the influence of friction is represented by

vxifi (vt - Lv),

where x is a small positive coefficient, in strictness the value of k is
also altered by the friction; but the alteration is of the second order as
regards the frictional forces, and may W omitted under the circumstances here
supposed. The energy of the waves per unit length at
any stage of degradation is proportional to the square of the amplitude,
and thus the whole energy on the positive side of the origin is to the
energy of so much of the weaves at their greatest value, i.e., at the
origin, as would be contained in the unit of length, as

or as

The energy transmitted through the origin in the
unit time is the same as the energy dissipated ; and, if the frictional
force acting on the element of mass m be
\todo[inline]{smudged}, where v is the velocity
of the element and h is constant, the energy dissipated in unit time is
\( h \sum m v^2\)
or
\(2 h T\), T being the kinetic energy. Thus, on the assumption
that the kinetic energy is half the whole energy, w e find that the
energy transmitted in the unit time to the greatest energy existing
in the unit length as h : ip.. It remains to find tin connection be
tween h and g..

For this purpose it will be convenient to regard cos (nt - kx) as the
real part of and to inquire how k is affected, when n is given
by the introduction of friction. Now the effect of friction is represented in
the differential equations of motion by the substitution of

-I- \(h dj/dt\)
in place of

or, since the whole motion is proportional to

by substituting \(-^2 + ihn\) for \(-n^2\).  Hence the introduction of
friction corresponds to an alteration of n from n to n- ih
(the square of h being neglected); and accordingly k is altered from k to
k-  The solution thus becomes g-i      or, when

the imaginary part is rejected, cos (ni - kx) ; so that

Li == h dk j dr i, and 4 : 2ft = dj i j dk. The ratio of the energy transmitted
in the unit time to the energy existing in the unit length is therefore
expressed by \(dn/dk\) or d (k V)jdk, as was to be proved.


It has often been noticed, in particular cases of progressive waves,
that the potential and kinetic energies are equal ; but I do not call to
mind any general treatment of the question. The theorem is not
usually true for the individual parts of the medium, but must be
understood to refer either to an integral number of wave-lengths, or to
a space so considerable that the outstanding fractional parts of waves
may be left out of account. As an example well adapted to give insight into the
question, I will take the case of a uniform stretched
circular membrane (Theory of Sound, § 200) vibrating with a given
number of nodal circles and diameters. The fundamental incites, are
not quite determinate in consequence of the symmetry, for any diameter may be
made nodal. In order to get rid of this indeterminateness, we may suppose the
membrane to carry a small load attached to
it anywhere except on a nodal circle. There are then two definite
fundamental modes, in one of which the node lies on a nodal diameter,
thus producing no effect, and in the other midway between nodal diameters,
where it produces a maximum effect (Theory of Sound, Article 208).
If vibrations of both modes are going on simultaneously, the
potential and kinetic energies of the whole motion may be calculated
by simple addition of those of the components. Let us now, supposing
the load to diminish without limit, imagine that the vibrations are of
equal amplitude and differ in phase by a quarter of a period. The
result is a progressive wave, whose potential and kinetic energies are
the sums of those of the stationary waves of which it is composed.
For the first component we have F, = cos and

for the second component, = E sin nt, T., - E cos at ; so that

Fj -I- Fj = -I- Tg = A, or the potential and kinetic energies of the
progressive wave are equal, being the same as the whole energy of
either of the components. The method of proof here employed appears
to be sufficiently general, though it is rather difficult to express it in
language which is appropriate to all kinds of waves.

 Aerial waves are an important exception.
END OF VOL. I.

People mentioned:
Airy, Arago, Bourget, Bravais, Chladni, Chree, Clebsch, Coddington, Colladon,
D'Alembert, Deschanel, Donkin, Duhem, Elsas, Ferrari, Froude, Germain, Giegsen,
Haweis, Hoppe, Jellet, Koenig, Kundt, Lamé, Lommel, Melde, Michell, Oberbeck,
Pockels, Preece, Riccati, Rijke, Routh, Savart, Scheibe, Scheibler, Sellmeier,
Senne, Somof, Stainbank, Strehlke, Stubbs, Sturm, Tait, Todhunter, Topler,
Wertheim, Young.
